\section{Fourier 基与 FFT}
\label{sec:03-Linear-Algebra/08-Complex-Matrices-and-Fourier-Structure/03}

给一段长度 \(n\) 的周期序列过一个滤波器：输出的每一项是输入在一个滑动窗口上的加权和。写成矩阵，它是一个循环矩阵——每一行都是上一行整体右移一格。换一组滤波系数，矩阵就换一个，这样的矩阵有无穷多个。

但它们全都由同一个东西生成。记循环移位

\[
(Px)_j=x_{j-1\bmod n},
\]

那么系数为 \(c_0,\ldots,c_{n-1}\) 的滤波器就是 \(C=\sum_{\ell}c_\ell P^\ell\)，一个关于 \(P\) 的多项式。多项式不会制造新的特征方向：只要找到 \(P\) 的特征向量，就一次性找到了所有循环卷积算子的特征向量。这个观察值得停一下——它意味着存在一组固定的坐标，把无穷多个滤波器同时变成对角矩阵。

\subsection{移位的特征方程只有一种解}

设 \(Pf=\lambda f\)。逐坐标写开是 \(f_{j-1}=\lambda f_j\)，所以

\[
f_j=\lambda^{-j}f_0 .
\]

关键在于下标是循环的：绕一圈回到起点，必须有 \(f_n=f_0\)，即 \(\lambda^{-n}=1\)。于是 \(\lambda\) 不能任选，只能是 \(n\) 次单位根。记 \(\omega=e^{2\pi i/n}\)，第 \(k\) 个解给出

\[
f^{(k)}_j=\frac1{\sqrt n}\,\omega^{kj},
\qquad
Pf^{(k)}=\omega^{-k}f^{(k)},
\qquad k=0,\ldots,n-1 .
\]

这组向量不是谁设计出来的，是周期性把它逼出来的。它们彼此正交，因为当 \(k\ne\ell\) 时 \(\sum_{j}\omega^{(k-\ell)j}=0\)——单位圆上一整圈等距点的和落在圆心。既然 \(P\) 是酉矩阵（移位不改变任何分量的模），上一节的结论适用：它可被酉对角化，而 \(f^{(0)},\ldots,f^{(n-1)}\) 正是那组标准正交特征基。

\begin{center}
\begin{tikzpicture}[x=1cm,y=1cm,font=\footnotesize,>=stealth]
  \draw[black!30] (0,0) circle (1.7);
  \draw[black!35] (-2.0,0) -- (2.0,0);
  \draw[black!35] (0,-2.0) -- (0,2.0);
  \foreach \k in {0,...,7} {
    \fill[blue!65!black] ({1.7*cos(45*\k)},{1.7*sin(45*\k)}) circle (2.3pt);
    \node at ({2.15*cos(45*\k)},{2.15*sin(45*\k)}) {\(\omega^{\k}\)};
  }
  \draw[black!60] (0.55,0) arc (0:45:0.55);
  \node[right,black!70] at (0.42,0.24) {\(2\pi/8\)};
  \node at (0,-2.75) {八次单位根};
  \begin{scope}[shift={(6.4,0)}]
    \draw[black!40] (-0.3,1.35) -- (5.0,1.35);
    \draw[blue!55!black,thin] plot[domain=0:4.55,samples=80] (\x,{1.35+cos(deg(2*pi*\x/5.2))});
    \foreach \j in {0,...,7} {
      \fill[blue!65!black] ({0.65*\j},{1.35+cos(45*\j)}) circle (2pt);
      \draw[blue!65!black] ({0.65*\j},1.35) -- ({0.65*\j},{1.35+cos(45*\j)});
    }
    \draw[black!40] (-0.3,-1.35) -- (5.0,-1.35);
    \draw[red!55!black,thin] plot[domain=0:4.55,samples=140] (\x,{-1.35+cos(deg(6*pi*\x/5.2))});
    \foreach \j in {0,...,7} {
      \fill[red!60!black] ({0.65*\j},{-1.35+cos(135*\j)}) circle (2pt);
      \draw[red!60!black] ({0.65*\j},-1.35) -- ({0.65*\j},{-1.35+cos(135*\j)});
    }
    \node at (2.2,-2.75) {\(k=1\)（上）与 \(k=3\)（下）的实部采样};
  \end{scope}
\end{tikzpicture}

\emph{左边八个点把圆周等分，它们的和为零，这就是不同频率相互正交的几何原因。右边是同一组点的另一种看法：第 \(k\) 个基向量沿圆周每次跨 \(k\) 格，投影到实轴便得到频率为 \(k\) 的余弦采样。\(k\) 越大，采样点绕圈越快，波形越密。Transformer 里的正弦位置编码用的就是这套``一列一个频率''的排布。}
\end{center}

\subsection{DFT 是一次精确的酉换基}

把这组特征向量的共轭转置排成矩阵，就得到单位化的离散 Fourier 变换（DFT）矩阵

\[
F_{kj}=\frac1{\sqrt n}\,\omega^{-kj},
\qquad 0\le k,j<n .
\]

信号的 Fourier 系数是 \(\widehat x=Fx\)。由于各列标准正交，\(F^{\conj}F=I\)，逆变换只是一次共轭转置：\(x=F^{\conj}\widehat x\)。所以换到频域是一次坐标改写，不是近似，也不损失信息；酉性顺带给出 Parseval 恒等式 \(\norm{x}_2^2=\norm{\widehat x}_2^2\)，时域和频域的总能量相同。数值上也因此安全——它不会放大误差。

取 \(n=4\)、\(x=(1,0,-1,0)^{\trans}\) 代进去，\(\widehat x_k=\tfrac12\bigl(1-(-1)^k\bigr)\)，即 \(\widehat x=(0,1,0,1)^{\trans}\)。时域上是一对正负交替的脉冲，频域上只剩一对频率，两边的平方范数都是 \(2\)。实信号的系数总满足 \(\widehat x_{n-k}=\overline{\widehat x_k}\)：正负频率成共轭对出现。这是实数条件在复指数坐标下的必然表现，看到频谱``对称''不要当成数据冗余或采集故障。

\subsection{卷积为什么退化成逐点相乘}

循环卷积

\[
(x*y)_j=\sum_{\ell=0}^{n-1}x_\ell\,y_{j-\ell\bmod n}
\]

固定 \(y\) 时就是开头那个循环矩阵 \(C=\sum_\ell y_\ell P^\ell\)。它与 \(P\) 共用特征向量，在第 \(k\) 个特征方向上的特征值是 \(\sum_\ell y_\ell\omega^{-k\ell}=\sqrt n\,\widehat y_k\)。换到 Fourier 坐标后，各频率互不往来，每个坐标只被自己的那个特征值乘一下：

\[
F(x*y)=\sqrt n\,(Fx)\odot(Fy),
\]

\(\odot\) 表示逐元素乘法。常数 \(\sqrt n\) 只是归一化约定的产物，若正变换不除 \(\sqrt n\)、逆变换除以 \(n\)，右边就没有它。不同教材的因子差异都出自这里，不是数学分歧。

\begin{center}
\begin{tikzpicture}[x=1cm,y=1cm,font=\footnotesize,>=stealth]
  \foreach \j in {0,...,6} {
    \draw[black!55] ({0.62*\j},0.62) rectangle ({0.62*\j+0.55},1.17);
  }
  \foreach \j in {2,3,4} {
    \draw[black!55,fill=blue!20] ({0.62*\j},0.62) rectangle ({0.62*\j+0.55},1.17);
  }
  \draw[black!55,fill=blue!8] (1.86,-0.65) rectangle (2.41,-0.1);
  \foreach \j in {2,3,4} {
    \draw[->,black!60] ({0.62*\j+0.28},0.55) -- (2.14,-0.02);
  }
  \foreach \j in {3,4,5} {
    \draw[->,black!30,dashed] ({0.62*\j+0.28},0.55) -- (2.76,-0.02);
  }
  \draw[black!35,dashed] (2.48,-0.65) rectangle (3.03,-0.1);
  \node at (2.0,-1.35) {时域：每个输出汇总一个窗口};
  \begin{scope}[shift={(6.6,0)}]
    \foreach \j in {0,...,5} {
      \draw[black!55] ({0.62*\j},1.05) rectangle ({0.62*\j+0.55},1.6);
      \draw[black!55] ({0.62*\j},0.25) rectangle ({0.62*\j+0.55},0.8);
      \draw[black!55,fill=blue!12] ({0.62*\j},-0.65) rectangle ({0.62*\j+0.55},-0.1);
      \draw[->,black!60] ({0.62*\j+0.28},0.2) -- ({0.62*\j+0.28},-0.05);
      \node[black!70] at ({0.62*\j+0.28},0.93) {\(\times\)};
    }
    \node[left] at (-0.1,1.32) {\(\widehat x\)};
    \node[left] at (-0.1,0.52) {\(\widehat y\)};
    \node at (1.7,-1.35) {频域：每个频率各算各的};
  \end{scope}
\end{tikzpicture}

\emph{左右是同一个运算的两种坐标。时域里每个输出格子都要去邻居那里取数，格子之间纠缠在一起；换到 Fourier 坐标后，纠缠全部消失，第 \(k\) 个频率的输出只依赖第 \(k\) 个频率的输入。频域滤波、谱方法解微分方程、以及下面的快速算法，都是在花式使用这一张右图。}
\end{center}

有一条隐含假设必须记住：上面的 \(*\) 是循环卷积，边界是绕回来的。要用 DFT 算普通的有限长线性卷积，得先给两个序列补足够多的零，让尾部绕回时撞不到有效数据。忘了补零，结果里就会出现信号尾巴叠到开头的环绕伪影。

\subsection{分治把 \texorpdfstring{\(n^2\) 降到 \(n\log n\)}{n 平方降到 n log n}}

直接按定义算，\(n\) 个系数每个要汇总 \(n\) 项，共 \(O(n^2)\) 次乘加。快速 Fourier 变换（FFT）算的是同一个 DFT，一个数值都不变，只是把重复利用的部分提出来。

机制是分治。把和式按下标的奇偶拆开，偶数项凑成一个长度 \(n/2\) 的 DFT，奇数项也凑成一个长度 \(n/2\) 的 DFT，两者用一个只依赖 \(k\) 的相位因子 \(\omega^{-k}\) 拼回去。这一步之所以成立，是因为 \(\omega^2\) 恰好是 \(n/2\) 次单位根——单位根的自相似性是全部加速的来源。于是

\[
T(n)=2T(n/2)+O(n),
\qquad T(n)=O(n\log n).
\]

\begin{center}
\begin{tikzpicture}[x=1cm,y=1cm,font=\footnotesize,>=stealth]
  \draw[black!55,fill=blue!10] (0,0) rectangle (7.2,0.6);
  \node at (3.6,0.3) {长度 \(n\)};
  \draw[black!55,fill=blue!10] (0,-1.1) rectangle (3.5,-0.5);
  \node at (1.75,-0.8) {偶数项 \(n/2\)};
  \draw[black!55,fill=blue!10] (3.7,-1.1) rectangle (7.2,-0.5);
  \node at (5.45,-0.8) {奇数项 \(n/2\)};
  \draw[->,black!60] (3.4,-0.05) -- (1.9,-0.45);
  \draw[->,black!60] (3.8,-0.05) -- (5.3,-0.45);
  \foreach \i in {0,1,2,3} {
    \draw[black!55,fill=blue!10] ({1.85*\i},-2.2) rectangle ({1.85*\i+1.65},-1.6);
    \node at ({1.85*\i+0.82},-1.9) {\(n/4\)};
  }
  \node[black!70] at (3.6,-2.7) {\(\vdots\)};
  \node[black!70] at (3.6,-3.2) {长度 \(1\) 的 DFT 就是它自己};
  \draw[->,black!45] (7.9,0.3) -- (7.9,-3.1);
  \node[right,black!70,align=left] at (8.0,-1.4) {共 \(\log_2 n\) 层\\每层拼接\\合计 \(O(n)\)};
\end{tikzpicture}

\emph{每往下一层，问题规模减半而个数加倍，所以每层的总工作量都是 \(O(n)\)；层数是 \(\log_2 n\)。两者相乘即 \(O(n\log n)\)。图上画的是最简单的二分版本，实际库对任意长度都有对应的分解策略，``FFT 只能用于 \(2^m\)''是对 radix-2 的误读。}
\end{center}

这条加速有直接的工程后果。长度 \(m\) 的卷积核直接滑动要 \(O(nm)\)，改成``两次正变换、一次逐点乘、一次逆变换''则是 \(O(n\log n)\)，与核的大小无关。所以大核卷积、长序列的自相关、多项式乘法都走 FFT；但常数不小，核只有三五个点时直接算反而更快，实际实现里存在一个明确的切换阈值。

\subsection{把线性图景推出去会在哪里断}

\hyperref[ch:126]{``卷积网络''}中的参数共享与这里的平移不变同源，一个不带偏置的循环卷积层确实能被 DFT 精确对角化。但网络还有通道混合、非线性激活、stride 与 pooling。非线性会把能量搬到新的频率上，下采样与有限边界破坏了理想的循环对称，所以``卷积在频域是乘法''能解释一层，不能对角化整个深度网络。

DFT 自身的假设同样需要盯住。它把观测窗口当作周期延拓，窗口两端接不上就会产生谱泄漏，加窗函数是缓解手段而不是修复。采样率不足时不同的连续频率会落在同一批离散样本上，这是混叠，信息已经在采样那一步丢了，增加 FFT 点数救不回来。零填充让频谱曲线画得更密，看上去分辨率变高，其实频率分辨率由观测时长决定，补零不产生新的观测。至于``信号在频域是稀疏的''，那是对信号结构的假设而非定理——一个尖锐的局部脉冲，它的 Fourier 系数会摊到几乎所有频率上。

这三节做的其实是同一件事：先确认标量域必须是 \(\C\)，再确认内积必须带共轭，然后让谱定理在这套语言里重新落地，最后发现平移结构会把特征基唯一确定下来。离散情形到此为止。把序列换成函数、把循环移位换成平移算子之后，同样的推理给出连续 Fourier 变换与谱方法，那是\hyperref[ch:068]{``Fourier 与谱方法''}的内容；而复特征值的模长与幅角怎样支配一个系统的振荡与衰减，会在\hyperref[ch:135]{``稳定性与动力系统''}里被反复调用。
